% ECONOMICS_PROBLEM: {"id": "C21-194", "title": "High-dimensional causal inference", "jel_code": "C21", "source_id": "OP-0104"}
% !TeX program = xelatex
\documentclass[11pt,a4paper]{article}
\usepackage[margin=23mm]{geometry}
\usepackage{fontspec}
\setmainfont{Latin Modern Roman}
\usepackage{amsmath,amssymb,mathtools}
\usepackage{xcolor,enumitem,booktabs,longtable,array,xurl,hyperref}
\definecolor{muted}{HTML}{53636C}
\hypersetup{colorlinks=true,urlcolor=blue,linkcolor=blue}
\setlength{\parindent}{0pt}
\setlength{\parskip}{5pt}
\newcommand{\R}{\mathbb R}
\newcommand{\E}{\mathbb E}
\newcommand{\N}{\mathbb N}
\newcommand{\ind}{\mathbf 1}
\newcommand{\Var}{\operatorname{Var}}
\newcommand{\Cov}{\operatorname{Cov}}
\newcommand{\supp}{\operatorname{supp}}
\DeclareMathOperator*{\argmin}{arg\,min}
\DeclareMathOperator*{\argmax}{arg\,max}
\newcommand{\entrymeta}[3]{{\sffamily\small\color{muted}\textbf{\texttt{#1}}\quad|\quad #2\par #3\par}}
\newcommand{\Problemref}[1]{\texttt{#1}}
\newcommand{\JELref}[2]{\texttt{#2} (JEL \texttt{#1})}
\begin{document}
\section*{High-dimensional causal inference}
\textbf{Problem ID:} \texttt{C21-194}\quad \textbf{JEL:} \texttt{C21}\par
% BEGIN ECONOMICS STATEMENT
\label{problem:OP-0104}
\label{atlas:prob:E4}

\noindent\textbf{Original identifier:} E4 (atlas item 104). \textbf{Field:} Econometrics and causal inference. \textbf{Classification:} editorial research family with established restricted solutions; current open status not certified.

\subsection*{Definitions and assumptions}

Let independent blocks $B_1,\ldots,B_m$ contain possibly dependent observations $(Y_i,D_i,X_i)$, with $X_i\in\mathbb R^{p_m}$ and $p_m$ potentially larger than $m$. Specify block-size growth and weighting before estimation. In a partially linear model $Y=\theta D+g(X)+\epsilon$ assume $\mathbb E[\epsilon\mid D,X]=0$, and let $l(X)=\mathbb E[Y\mid X]$, $q(X)=\mathbb E[D\mid X]$. Require $\mathbb E[(D-q(X))^2]>0$. The score
\[
\psi(W;\theta,l,q)=(D-q(X))\{Y-l(X)-\theta(D-q(X))\},
\qquad W=(Y,D,X),
\]
has zero expectation at the true parameter. A causal interpretation requires structural invariance, for example $Y(d)=\theta d+g(X)+\epsilon$, and suitable assignment restrictions.

\subsection*{Formal research question}

Characterize model classes permitting uniform inference for $\theta$ with block-level cross-fitting, growing covariate dimension and distribution drift. A baseline sufficient rate condition for independent effective observations is
\[
\|\widehat q-q\|_2
\bigl(\|\widehat l-l\|_2+\|\widehat q-q\|_2\bigr)
=o_P(m^{-1/2}),
\]
along with suitable moments, nondegenerate score variance and consistent variance estimation; $\|f\|_2=(\mathbb E[f(X)^2])^{1/2}$. Determine appropriate replacements when block dependence reduces effective information or nuisance functions differ by environment. Seek confidence sets $C_m$ satisfying uniform $1-\alpha$ coverage, for $\alpha\in(0,1)$, over a stated approximately sparse or otherwise learnable class, and derive lower bounds when nuisance complexity, weak residual variation or drift prevents the required rates.

\subsection*{Interpretation and scope}

Folds must keep dependent blocks intact; splitting individual observations across folds need not make training and evaluation independent. The number of predictors alone is not the difficulty: unrestricted confounding cannot be learned from finite data without complexity restrictions. Approximate sparsity needs a quantified approximation error; prediction accuracy and causal identification are separate requirements. The partially linear coefficient is a constant causal effect only under its stated structural model. If treatment effects are heterogeneous, specify another estimand and an appropriate orthogonal score instead of relabeling this slope as an ATE.

\subsection*{Source audit and status}

[S1] establishes post-selection inference under approximate sparsity; [S2] gives double/debiased ML theory under explicit regularity conditions. The undated Athey--Imbens--Wager citation is completed with the verified 2018 approximate-residual-balancing article [S3]; intended-source certainty is limited because the input supplied no title. Standard independent-data high-dimensional inference has established solutions. The dependent and drifting regime above is a scoped editorial frontier, not a published conjecture certified open today.

\subsection*{References}

\begin{enumerate}[label={[\textnormal{S}\arabic*]},leftmargin=*]

\item Alexandre Belloni, Victor Chernozhukov and Christian Hansen (2014). \emph{Inference on Treatment Effects after Selection among High-Dimensional Controls}. Review of Economic Studies 81(2), 608–650. \url{https://doi.org/10.1093/restud/rdt044}. \textbf{Locator:} Publisher or institutional abstract and bibliographic record. \textbf{Support and limits:} Post-double-selection inference under approximate sparsity and regularity; not validity for arbitrary high-dimensional confounding.

\item Victor Chernozhukov, Denis Chetverikov, Mert Demirer, Esther Duflo, Christian Hansen, Whitney Newey and James Robins (2018). \emph{Double/debiased machine learning for treatment and structural parameters}. Econometrics Journal 21(1), C1–C68. \url{https://doi.org/10.1111/ectj.12097}. \textbf{Locator:} Publisher or institutional abstract and bibliographic record. \textbf{Support and limits:} Orthogonal scores and cross-fitting under nuisance-rate and sampling restrictions; not automatic protection against dependence or misspecification.

\item Susan Athey, Guido W. Imbens and Stefan Wager (2018). \emph{Approximate Residual Balancing: Debiased Inference of Average Treatment Effects in High Dimensions}. Journal of the Royal Statistical Society Series B 80(4), 597–623. \url{https://doi.org/10.1111/rssb.12268}. \textbf{Locator:} Publisher or institutional abstract and bibliographic record. \textbf{Support and limits:} Verified three-author candidate; the input omitted year and title. Selected explicitly as a relevant completion, not proof of intended citation.

\end{enumerate}

\subsection*{Common conventions: Source mathematical conventions}

Unless an entry states otherwise, agent and firm sets are finite. State and action spaces are measurable, strategies and outcome maps are measurable, probability laws are specified, and expectations exist for the targets under discussion. Infinite-horizon discount factors lie in $(0,1)$ and discounted objectives are finite. Norms, tolerances and complexity input sizes must be specified for computational comparisons. Constants or primitives that appear locally are inputs, not empirical facts asserted by this catalogue.

\textbf{Identification.} A statistical model $m\in\mathcal M$ implies an observable law $P_m$ and target $\theta(m)$. At law $P$, its identified set is
\[
\Theta(P;\mathcal M)=\{\theta(m):m\in\mathcal M,\ P_m=P\}.
\]
Point identification holds when this set is a singleton, or equivalently when $P_{m_1}=P_{m_2}$ implies $\theta(m_1)=\theta(m_2)$ for all compatible models. Sharp bounds mean bounds attained, or approached when the boundary is not attained, by compatible models. Writing this set is a definition, not a solution: the substantive task is to establish informative restrictions and derive or estimate the set. An unrestricted-confounding model may give only trivial bounds.

\textbf{Inference.} Identification, estimability and valid uncertainty quantification are separate questions. Fix $\alpha\in(0,1)$. A confidence procedure $C_n$ over a declared law class $\mathcal P$ has asymptotic uniform coverage at level $1-\alpha$ when
\[
\liminf_{n\to\infty}\inf_{P\in\mathcal P}\Pr_P\{\theta(P)\in C_n\}\geq1-\alpha.
\]
For set-identified targets the coverage event must be defined separately, for example coverage of the full identified set. If an entry uses identification-indexed models instead of uniquely indexed laws, the same coverage convention is applied over those models. Pointwise limits do not establish uniform validity, and asymptotic validity does not guarantee finite-sample precision. Sample sizes, dependence structures and weak-identification sequences are part of each application.

\textbf{Counterfactuals.} $Y_i(a)$ is a potential outcome under a specified intervention, including its timing, implementation and versions. It is not $Y_i$ conditional on voluntarily choosing $a$. A full intervention vector permits interference; shorthand scalar treatment notation does not silently impose no interference. Assignment, selection, attrition, anticipation and measurement must be specified. A claim about an instrument's local effect is not automatically a population policy effect.

\textbf{Economic equilibrium.} A counterfactual model includes best responses, constraints, feasibility, market clearing, state transitions and expectations consistent with the stated information. When several equilibria exist, either a declared selector is an input or the target is set-valued. Comparisons across policy regimes must use the stated selection convention. Reduced-form relationships are not presumed invariant after an equilibrium-changing intervention.

\textbf{Optimization and welfare.} Optimization is over the stated feasible set. If attainment is not established, $\sup$ or $\inf$ and $\varepsilon$-optimal choices replace an asserted maximizer or minimizer. For example, with value $V=\sup_{a\in\mathcal A}W(a)<\infty$, an $\varepsilon$-optimal set is $\{a\in\mathcal A:W(a)\geq V-\varepsilon\}$ for $\varepsilon>0$. Benefits, costs, risk and health or environmental outcomes can be aggregated only using declared units and conversion weights. Interpersonal utility weights, the treatment of fiscal transfers and foreign profits, discounting, and the behavioral welfare criterion are explicit inputs. A comparative-static derivative additionally requires existence and appropriate differentiability of the selected counterfactual.

\textbf{Sources.} Local labels [S1], [S2], and so forth restart in every question. Locators identify the relevant abstract, model, section or result at the level actually checked; a bibliographic or abstract-level verification is not represented as inspection of an entire proof. Working papers are distinguished from later publications and changing versions. Source summaries are paraphrased. All URL checks and editorial formulations refer to the audit date stated above.
% END ECONOMICS STATEMENT
\end{document}
